\textbf{Lösung zu Übung 6, Aufgabe 9}

\begin{loesung}
Inverse Modulo 11
\begin{center}
\begin{tabular}{c|rrrrrrrrrr}
$a$      & 1 & 2  & 3 & 4 & 5  & $-5$ & $-4$ & $-3$ & $-2$ & $-1$\\
\midrule
$a^{-1}$ & 1 & $-5$ & 4 & 3 & $-2$ & 2 & $-3$ & $-4$ & 5 & $-1$\\
\end{tabular}
\end{center}

\begin{align*}
\mathbf{r}_1 &= 1\,1\,5\,1\,1\,1\,1\,1\,0\,3\\
\mathbf{r}_2 &= 1\,1\,5\,6\,1\,1\,1\,1\,0\,3\\
\mathbf{r}_3 &= 1\,1\,1\,1\,1\,1\,1\,1\,0\,3
\end{align*}

\[
H = \begin{pmatrix}
1 & 1 & 1 & 1 & 1 & 1 & 1 & 1 & 1 & 1\\
1 & 2 & 3 & 4 & 5 & 6 & 7 & 8 & 9 & 10
\end{pmatrix}
\equiv
\begin{pmatrix}
1 & 1 & 1 & 1 & 1 & 1 & 1 & 1 & 1 & 1\\
1 & 2 & 3 & 4 & 5 & -5 & -4 & -3 & -2 & -1
\end{pmatrix}
\]

$\mathbf{e} = \mathbf{r} - \mathbf{c} = (\,0 \dots m \dots 0\,)$, \quad
$S(\mathbf{r}) = \begin{pmatrix} s_0\\ s_1 \end{pmatrix} = \begin{pmatrix} 1\\ x \end{pmatrix} y$, \quad
$\mathbf{c} = \mathbf{r} - \mathbf{e}$

\[
S(\mathbf{r}_1) = \begin{pmatrix} 15\\ 4 \cdot 3 \end{pmatrix} = \begin{pmatrix} 4\\ 4 \cdot 3 \end{pmatrix} = \begin{pmatrix} 1\\ 3 \end{pmatrix} 4
\qquad\Rightarrow\ x = 3,\ y = 4
\]
$\Rightarrow \mathbf{c}_1 = 1\,1\,1\,1\,1\,1\,1\,1\,0\,3 = \mathbf{r}_3$

\[
S(\mathbf{r}_2) = \begin{pmatrix} 9\\ 1 + 5 \cdot 4 \end{pmatrix} = \begin{pmatrix} -2\\ -1 \end{pmatrix} = \begin{pmatrix} 1\\ 6 \end{pmatrix} (-2)
\qquad\Rightarrow\ x = 6,\ y = -2
\]
$\Rightarrow \mathbf{c}_2 = 1\,1\,5\,6\,1\,3\,1\,1\,0\,3$

\[
S(\mathbf{r}_3) = \begin{pmatrix} 0\\ 0 \end{pmatrix} \qquad\Rightarrow \mathbf{c}_3 = \mathbf{r}_3
\]
\end{loesung}
